% mycsrf 'for beeing included' snippet template
%
% (c) Karsten Reincke, Frankfurt a.M. 2012, ff.
%
% This text is licensed under the Creative Commons Attribution 3.0 Germany
% License (http://creativecommons.org/licenses/by/3.0/de/): Feel free to share
% (to copy, distribute and transmit) or to remix (to adapt) it, if you respect
% how you must attribute the work in the manner specified by the author(s):
% \newline
% In an internet based reuse please link the reused parts to mycsrf.fodina.de
% and mention the original author Karsten Reincke in a suitable manner. In a
% paper-like reuse please insert a short hint to mycsrf.fodina.de and to the
% original author, Karsten Reincke, into your preface. For normal quotations
% please use the scientific standard to cite
%


Wenn wir nun die bisher gewonnenen Informationen verknüpfen, gewinnen wir folgende Übersicht:


\begin{tikzpicture}
\footnotesize

\tikzstyle{every node}=[trapezium, draw, minimum width=1cm,trapezium left angle=70, trapezium right angle=110]

% BOTTOM UP DESCRIPTION of the component IO structure:
% THE FINALLY GENERATED PDF OUTPOUT ->
\file{\pdfBgColor}{black}{PDF}{5.5}{0}{.pdf}

% -> CREATED BY THE LATEX IDEs ->
\backend{$\star\star\star$}{\abcColor}{LTXABC}{1}{2}{\LaTeX\,+\,ABC}
\backend{$\star\star\star\star$}{\mtexColor}{LTXMTEX}{5.5}{2}{\LaTeX\,+\,MusiX\TeX}
\backend{$\star\star\star\star\star$}{\lyColor}{LTXLIPO}{10.4}{2}{\LaTeX\,+\,LilyPond}

% -> which take the CONVERTER OUTPUT ->
\file{\abcBgColor}{\abcColor}{CO-ABC}{1}{3}{.abc}
\file{\mtexBgColor}{\mtexColor}{CO-MTEX}{5.5}{3}{.mtex}
\file{\lyBgColor}{Blue}{CO-LY}{10.4}{3}{.ly}

% written by the CONVERTERs.
\converter{\mxmlInColor}{\abcOutColor}{XML2ABC}{0}{7}{xml2abc}
\converter{\mxmlInColor}{\pmxOutColor}{MXML2PMX}{3.6}{8.2}{mxml2pmx}
\converter{\mxmlInColor}{\lyOutColor}{XML2LY}{5}{10.5}{xml2ly}
\converter{\mxmlInColor}{\lyOutColor}{MXML2LY}{5}{12}{mxml2ly}

\converter{\abcInColor}{\mtexOutColor}{ABC2MTEX}{3}{4}{abc2mtex}
\converter{\abcInColor}{\lyOutColor}{ABC2LY}{5}{11.2}{abc2ly}

\converter{\lyInColor}{\abcOutColor}{LY2ABC}{5.6}{8}{ly2abc}
\converter{\pmxInColor}{\mtexOutColor}{PMXAB}{5.5}{4}{pmxab}


% These converters takte the FRONTEND OUTPUT ->

\file{\abcBgColor}{\abcColor}{FO-ABC}{0}{12}{.abc}
\file{\mxmlBgColor}{\mxmlColor}{FO-MXML}{1.5}{12}{.mxml}
\file{\pmxBgColor}{Brown}{FO-PMX}{5.5}{5.6}{.pmx}
\file{\lyBgColor}{Blue}{FO-LY}{7.2}{10.5}{.ly}

% created by the FRONTENDS ->

\frontend{\abcBgColor}{Black}{EASYABC}{0}{13.8}{EasyABC}{$\star\star\star\star$};
\frontend{\lyBgColor}{Black}{ROSEGARDEN}{2.2}{14.8}{Rosegarden}{$\star\star\star$};
\frontend{\multiBgColor}{Black}{MUSESCORE}{4.5}{14.8}{MuseScore}{$\star\star\star\star\star$};
\frontend{\lyBgColor}{Black}{DENEMO}{6.5}{14.8}{Denemo}{$\star\star\star$};
\frontend{\lyBgColor}{Black}{CANORUS}{8.3}{14.8}{Canorus}{$\star\star\star$};

\frontend{\multiBgColor}{Black}{EDITOR}{7.7}{7.2}{Text-Editor}{$\star$};
\frontend{\lyBgColor}{Black}{ELYSIUM}{9.8}{7.6}{Elysium}{$\star\star\star\star$};
\frontend{\lyBgColor}{Black}{FRESCOBALDI}{11.9}{8}{Frescobaldi}{$\star\star\star\star\star$};

% which are used by the Composer:

%
\node[rectangle, draw=none, inner sep=0pt] (Composer) at (8.3,16.5) 
{\includegraphics{./logos/decider-220x320.eps}};
  
% TOP DOWN linkings of the component IO structure:  
  
% (1) The composer uses the frontends.
\draw[->] (Composer) to [out=180,in=70] (EASYABC);  
\draw[->] (Composer) to [out=180,in=50] (ROSEGARDEN); 
\draw[->] (Composer) to [out=180,in=70] (MUSESCORE);  
\draw[->] (Composer) to [out=180,in=80] (DENEMO);  
\draw[->] (Composer) to [out=270,in=90] (CANORUS);  
\draw[->] (Composer) to [out=0,in=90] (ELYSIUM);  
\draw[->] (Composer) to [out=0,in=90] (FRESCOBALDI);  
\draw[->] (Composer) to [out=0,in=90] (EDITOR);  

 
% (2.A) The frontends write the frontend output  
\draw[->,\abcColor] (EASYABC) to [out=270,in=90] (FO-ABC); 
\draw[->,\abcColor] (EASYABC) to [out=240,in=140] (CO-ABC); 
\draw[->,\mxmlColor] (EASYABC) to [out=300,in=90] (FO-MXML); 

\draw[->,\lyColor] (ROSEGARDEN) to [out=300,in=90] (FO-LY); 
\draw[->,\mxmlColor] (ROSEGARDEN) to [out=240,in=90] (FO-MXML);

\draw[->,\mxmlColor] (MUSESCORE) to [out=240,in=90] (FO-MXML); 

\draw[->,\mxmlColor] (CANORUS) to [out=230,in=70] (FO-MXML); 
\draw[->,\lyColor] (CANORUS) to [out=270,in=90] (FO-LY); 

\draw[->,\mxmlColor] (DENEMO) to [out=240,in=90] (FO-MXML); 
\draw[->,\lyColor] (DENEMO) to [out=270,in=90] (FO-LY); 


\draw[->,\abcColor] (EDITOR) to [out=270,in=80] (CO-ABC);
\draw[->,\mtexColor] (EDITOR) to [out=270,in=0] (CO-MTEX);
\draw[->,\pmxColor] (EDITOR) to [out=270,in=90] (FO-PMX);

\draw[->,\lyColor] (EDITOR) to [out=270,in=90] (CO-LY);
\draw[->,\lyColor] (ELYSIUM) to [out=270,in=90] (CO-LY); 
\draw[->,\lyColor] (FRESCOBALDI) to [out=270,in=90] (CO-LY); 

% (3) The frontend output becomes input of the converter

\draw[->,\abcColor] (FO-ABC) to [out=270,in=90] (ABC2MTEX);
\draw[->,\abcColor] (FO-ABC) to [out=270,in=180] (ABC2LY);

\draw[->,\mxmlColor] (FO-MXML) to [out=270,in=90] (XML2ABC);;
\draw[->,\mxmlColor] (FO-MXML) to [out=270,in=180] (XML2LY);
\draw[->,\mxmlColor] (FO-MXML) to [out=270,in=180] (MXML2LY);
\draw[->,\mxmlColor] (FO-MXML) to [out=270,in=90] (MXML2PMX);

\draw[->,\lyColor] (FO-LY) to [out=270,in=90] (LY2ABC); 
%\draw[->,\lyColor] (FO-LY) to [out=270,in=90] (EDITOR); 
\draw[->,\lyColor] (FO-LY) to [out=270,in=90] (ELYSIUM); 
\draw[->,\lyColor] (FO-LY) to [out=270,in=90] (FRESCOBALDI); 

% (3) The converters (re)write their output
% (3.A) as input for the textual frontends
\draw[->,\lyColor] (ABC2LY) to [out=0,in=90] (FO-LY);
\draw[->,\lyColor] (MXML2LY) to [out=0,in=90] (FO-LY); 
\draw[->,\lyColor] (XML2LY) to [out=0,in=180] (FO-LY); 

% (3.B) as input for the backends

\draw[->,\abcColor] (LY2ABC) to [out=270,in=90] (CO-ABC);
\draw[->,\abcColor] (XML2ABC) to [out=270,in=90] (CO-ABC);

\draw[->,\pmxColor] (MXML2PMX) to [out=270,in=90] (FO-PMX);
\draw[->,\mtexColor] (ABC2MTEX) to [out=270,in=180] (CO-MTEX);

% special case MusiXTeX
\draw[->,\pmxColor] (FO-PMX) to [out=270,in=90] (PMXAB);
\draw[->,\mtexColor] (PMXAB) to [out=270,in=90] (CO-MTEX);

% 4) The backends take the (converter) output as input
\draw[->,\abcColor] (CO-ABC) to [out=270,in=90] (LTXABC);
\draw[->,\mtexColor] (CO-MTEX) to [out=270,in=90] (LTXMTEX);
\draw[->,\lyColor] (CO-LY) to [out=270,in=90] (LTXLIPO);

% 5) end write the final PDF

\draw[->,Black] (LTXABC) to [out=270,in=90] (PDF);
\draw[->,Black] (LTXMTEX) to [out=270,in=90] (PDF);
\draw[->,Black] (LTXLIPO) to [out=270,in=90] (PDF);
\end{tikzpicture}

Wir kennen demnach vier echte graphische Frontends -- \acc{Rosegarden},
\acc{MuseScore}, \acc{Denemo} und \acc{Canorus} --, drei semi-graphische
Frontends -- \acc{EasyABC}, \acc{Elysium} und \acc{Frescobaldi} --,  drei
Backendsysteme -- \acc{\LaTeX\,+\,ABC}, \acc{\LaTeX\,+\,MusiX\TeX} und
\acc{\LaTeX\,+\,LilyPond} -- und eine Reihe von Konvertern. Unglücklicherweise
haben wir lernen müssen, dass es momentan keinen graphischen oder
semi-graphischen Editor für \acc{MusiX\TeX} gibt. Deshalb müssen wir ersatzweise
den reinen Texteditor in unsere Kollektion von Frontends aufnehmen.\footnote{Um
die Analyse nicht zu verfälschen, geben wir dem Text-Edtor nur einen Stern.}
Außerdem haben wir festgestellt, dass keines der echten graphischen Frontends
eine optimale Dokumentation der Harmonieanalyse bietet.
Und schließlich hatte sich schon vorher ergeben, dass das
\acc{\LaTeX-ABC-Backend} nur die Integration begrenzt komplexer Analysesymbole
anbietet, während die Kombinationen von \LaTeX\ und MusiX\TeX\ bzw. \LaTeX und
LilyPond jeweils gut mit strukturell adäquaten Analysesymbolen umgehen können.
Allerdings sind diese Kennzeichnungen bei der \LaTeX-MusiX\TeX-Variante z.Zt.
visuell ansprechender, als bei der \LaTeX-LilyPond-Variante.

Noch ist unsere 'Landkarte' allerdings bloß eine theoretische Skizze: die Wege
ergaben sich rein aus dem 'Versprechen' der Tools, bestimmte Formate zu
schreiben und zu lesen. Über die Praxis sagt das nichts. Die werden wir Weg für
Weg evaluieren müssen. Gleichwohl sollten wir die Güte eines Weges dann auch
systematisch vergleichbar in einer Zahl repräsentieren können. Glücklicherweise
bietet sich dafür unmittelbar ein Bewertungsverfahren an:

\begin{enumerate}
  \item Wir haben die Frontend- und Backendsysteme bereits gemäß ihrer immanenten
  Fähigkeiten klassifiziert. Und jeder 'Weg' beginnt mit einem Frontend und endet
  an einem Backend. Also sollte die Summe der Frontend- und Backend-Sterne etwas
  über die Qualität des Gesamtsystems aussagen.
  \item Eine erhöhte Komplexität des Weges bildet jedoch einen Nachteil: je mehr
  Konvertierungsschritte nötig sind, desto größer ist der Aufwand für den
  Musikwissenschaftler. Es ist also angemessen, wenn wir von der Summe der
  Sterne die Anzahl der Konvertierungsschritte wieder abzuziehen: je direkter
  der Weg, desto besser.\footnote{Da aber zwischen einem Frontend und einem
  Backend immer mindestens ein Konvertierungsschritt liegt, der sich in der
  geschriebenen Outputdatei des Frontends manifestiert, die zum Input des
  Backends wird, brauchen wir letztlich von Summer der Sterne nur die jeweiligen
  Konvertierungsschritte - 1 zu subtrahieren.}
  \item In den Bewertung der Alternativen muss außerdem die Qualität des
  generierten Ergebnisses eingerechnet werden. Wir wollen sagen, dass es
  \acc{exzellent}, \acc{akzeptabel}, \acc{grenzwertig} oder \acc{inakzeptabel}
  sein kann.
  \item Ferner wird nicht nur die Anzahl der Konvertierungsschritte die
  Nutzbarkeit eines Weges beeinflusses, sondern auch die technische
  Umsetzbarkeit: Was vom Format her -- sozusagen auf dem Papier --
  \acc{problemlos} verknüpfbar ist, kann immer noch in technische Probleme
  laufen, die das Verfahren \acc{stolpern} lassen (und leichte manuelle
  Eingegriffe nötig machen), es \acc{beschränken} (und starke manuelle Eingriffe
  erfordern) oder gar \acc{unmöglich} machen.
\end{enumerate}
 
Die Qualität des Ergebnisses und die Nutzbarkeit des Weges werden zuletzt
empirisch zu ermitteln sein. Wir werden die beiden Werte allerdings nach einem
festen Schema in einen einzigen Wert der \acc{Prozessgüte} verrechnen: Die Qualität
eines 'Weges', der problemlos ein exzellentes Ergebnis erzeugt, ist so gut, wie
die Qualität des beteiligten Frontends und Backends abzgl. der Komplexität des
Weges. Das wird mit einer gewichtenden Prozessgüte von 1 erreicht. Umgekehrt
sind die Wege nutzlos, die inakzeptable Ergebnisse erzeugen oder nicht umsetzbar
sind. Dem entspricht dann eine faktorelle Prozessgüte von 0. Alle anderen
Abstufungen ergeben sich aus folgender Tabelle:

\begin{center}
\renewcommand{\arraystretch}{1.5}
\begin{tabulary}{14cm}{|C|R||C|C|C|C||}
\cline{3-6}
 \multicolumn{2}{c||}{\ } & \multicolumn{4}{c||}{Technische Ausführbarkeit} \\
\cline{3-6}
 \multicolumn{2}{c||}{\ } & problemlos & stolpernd & beschränkt & unmöglich \\
\hline
\multirow{4}{*}{\rotatebox{90}{Ergebnis}} & exzel\-lent & \centering{1} & 0.8 & 0.6 & 0 \\
\cline{2-6}
 & akzep\-tabel & 0.8 & 0.6 & 0.4 & 0 \\
\cline{2-6}
 & grenz\-wertig & 0.6 & 0.4 & 0.2 & 0 \\
\cline{2-6}
 & inakzep\-tabel & 0 & 0 & 0 & 0 \\
\hline
\hline
\end{tabulary}
\renewcommand{\arraystretch}{1}
\end{center}

Damit lässt sich die Qualität eines Weges nach einer einfachen Methode
rechnerisch vergleichen: 

\label{ProzessReife}
\begin{description}
  \item[D0] Sei PG die nach obiger Tabelle empirisch ermitteltete Prozessgüte.
  \item[D1] Seien $stars(BE)$ und $stars(FE)$ die Anzahl der bis zu maximal 5
  Sterne, mit denen wir die Frontends und Backends klassifiziert haben.
  \item[D2] Sei $\Sigma$ die Anzahl der zwischenzuschaltenden Konverter,
  um den Output des Frontends dem Backend als Input vorlegen zu können.
  \item[D3] Sei $PR$ die Prozessreife, die nach der Formel $(stars(FE) +
  stars(BE) - \Sigma) * PG$ berechnet wird.
  \item[L1] Die maximale Prozessreife ist 10. Beweis: Frontends und Backends
  haben maximal 5 Sterne (\ra D1). Der direkte Weg hat 0 Konvertierungsschritte
  (\ra D2). Und die beste Prozessgüte hat den Wert 1 (\ra D1). Also gilt
  $max(PR) = (5 + 5 - 0) * 1 = 10$
  \item[D4] Sei $NPR$ die normierte Prozessreife, berechnet gemäß der Formel
  $PR/max(PR)$.
\end{description}

Umgangssprachlich würde man sagen, dass die Systemqualität eines 'Weges' in \% der
maximal möglichen Prozessreife angegeben wird. 

Bei einigen 'Wegen' resp. Gesamtsystemen können wir die Prozessgüte direkt
angeben: Wir wissen bereits, dass \acc{Frescobaldi} und \acc{Elysium} zur
Visualisierung des Inputs das \acc{LilyPond-Backend} benutzen. Und wir haben
diese Komponenten bereits -- zusammen mit unser kleinen Zusatzbibliothek --
erfolgreich bei unser Evaluierung der Backendsystem benutzt. Wir dürfen hier
also direkt eine Prozessgüte von 1 ansetzen. Dasselbe gilt für den -- sozusagen
nachträglich eingeführten -- Texteditor: Wir haben den Input zur Evaluierung des
\acc{\LaTeX-ABC-Backends} ebenso manuell mit dem \acc{Eclipse-Editor}
geschrieben, wie den für das \acc{\LaTeX-MusiX\TeX-Backend} und den für das
\acc{\LaTeX-LilyPond-Backend}. Mithin ergibt sich auch hier Prozessgüte von 1.
Und wir hatten die Schwierigkeiten bei der Verwendung von \acc{PMX} mit
\acc{\LaTeX\ + MusiX\TeX} ausführlich dokumentiert: Hier konnte zuletzt zwar ein
exzellentes Ergebnis erzeugt werden, allerdings nur um den Preis starker
manueller Eingriffe. Das ergibt per Definitionem eine Prozessgüte von $0.6$.

Für diese soeben aufgezählten Varianten können wir mithin auch die normierte
Prozessreife direkt berechnen. Die Prozessgüte der anderen Wegen werden wir
schrittweise evaluieren müssen, um danach die entsprechende Prozessreife
berechnen zu können.

Bringt man all dies zusammen, ergibt sich folgendes vorläufiges Bild:

\begin{footnotesize}

\begin{longtable}{|c||c|c|c|c|c|c||c||c||}

% was soll am Anfang oben scheinen:
\hline
  \multirow{2}{*}{\rotatebox{90}{Weg}} & \multicolumn{2}{|c}{VON} & \multicolumn{2}{|c}{ÜBER} & 
  \multicolumn{2}{|c||}{NACH} & \multicolumn{2}{c||}{WERTUNG} \\
\cline{2-9}
 & Frontend & $\bigstar$ & Konvertierung & $\Sigma$  & \LaTeX\ + & $\bigstar$ & 
  PG\tnote{a} & NPR\tnote{b}\\
\hline
\hline
\endfirsthead

% was soll bei jedem Seitenumnruch oben scheinen:
\multicolumn{4}{c}%
{\tablename\ \thetable\ } \\
\hline
 \multirow{2}{*}{\rotatebox{90}{Weg}} & \multicolumn{2}{|c}{VON} & \multicolumn{2}{|c}{ÜBER} & 
  \multicolumn{2}{|c||}{NACH} & \multicolumn{2}{c||}{WERTUNG} \\
\cline{2-9}
  & Frontend & $\bigstar$ & Konvertierung & $\Sigma$  & \LaTeX\ + & $\bigstar$ & 
  PG\tnote{a} & NPR\tnote{b}\\
\hline
\hline
\endfirsthead

% was soll bei jedem Seitenumnruch unter scheinen:
\hline \multicolumn{4}{r}{\textit{Continued on next page}} \\
\endfoot

% was am Ende erscheinen:
\hline
\endlastfoot

01 & Easy\-ABC & 4 & \ra\ & 0 &  ABC & 3 & &  \\
\hline
02 &  Easy\-ABC & 4 & \ra\ xml2abc \ra\  & 1 &  ABC & 3 & & \\
\hline
04 &  Easy\-ABC & 4 & \ra\ abc2ly \ra\ ly2abc \ra\ & 2 &  ABC & 3 & & \\
\hline
05 &  Easy\-ABC & 4 & \ra mxml2ly \ra\  ly2abc \ra\ & 2 &  ABC & 3 & & \\
\hline
06 &  Easy\-ABC & 4 & \ra xml2ly \ra\  ly2abc \ra\ & 2 &  ABC & 3 & & \\
\hline
07 &  Rose\-garden & 3 & \ra\ xml2abc \ra\ & 1 & ABC & 3 & &  \\
\hline 
09 & Rose\-garden & 3 & \ra\ mxml2ly \ra\ ly2abc \ra\ & 2 &  ABC & 3 &  &  \\
\hline
10 & Rose\-garden & 3 & \ra\ xml2ly \ra\ ly2abc \ra\ & 2 &  ABC & 3 &  &  \\
\hline
11 &  Muse\-Score & 5 & \ra\ xml2abc \ra\ & 1 & ABC & 3 &  &  \\
\hline
13 &  Muse\-Score & 5 & \ra\ mxml2ly \ra\ ly2abc \ra\ & 2 &  ABC & 3 &  &  \\
\hline
14 &  Muse\-Score & 5 & \ra\ xml2ly \ra\ ly2abc \ra\ & 2 &  ABC & 3 &  &  \\
\hline
15 &  Denemo & 3 & \ra\ xml2abc \ra\ & 1 & ABC & 3 &  &  \\
\hline
17 &  Denemo & 3 & \ra\ mxml2ly \ra\ ly2abc \ra\ & 2 &  ABC & 3 &  &  \\
\hline
18 &  Denemo & 3 & \ra\ xml2ly \ra\ ly2abc \ra\ & 2 &  ABC & 3 &  &  \\
\hline
19 &  Canorus & 3 & \ra\ xml2abc \ra\ & 1 & ABC & 3 &  &  \\
\hline
21 &  Canorus & 3 & \ra\ mxml2ly \ra\ ly2abc \ra\ & 2 &  ABC & 3 &  &  \\
\hline
22 &  Canorus & 3 & \ra\ xml2ly \ra\ ly2abc \ra\ & 2 &  ABC & 3 &  &  \\
\hline
23 &  Texteditor & 1 & \ra\ & 0 &  ABC & 3 & 1.0 & 0.4 \\
\hline
\hline
24 & EasyABC & 4 & \ra\ abc2mtex \ra\ & 1 &  Musix\TeX\ & 4 &  &  \\
\hline
25 & EasyABC & 4 & \ra\ mxml2pmx \ra\ pmxab \ra\ & 2 &  Musix\TeX\ & 4 & &  \\
\hline
26 & Rosegarden & 3 & \ra\ mxml2pmx \ra\ pmxab \ra\ & 2 &  Musix\TeX\ & 4 &  & \\
\hline
27 & MuseScore & 5 & \ra\ mxml2pmx \ra\ pmxab \ra\ & 2 &  Musix\TeX\ & 4 &  & \\
\hline
28 & Denemo & 3 & \ra\ mxml2pmx \ra\ pmxab \ra\ & 2 &  Musix\TeX\ & 4 &  &  \\
\hline
29 & Canorus & 3 & \ra\ mxml2pmx \ra\ pmxab \ra\ & 2 &  Musix\TeX\ & 4 &  &  \\
\hline
30 & Texteditor & 1 & \ra\ & 0 &  Musix\TeX\ & 4 & 1.0 & 0.5 \\
\hline
31 & Texteditor & 1 & \ra\ pmxab \ra\ & 1 &  Musix\TeX\ & 4 & 0.6 & 0.24 \\
\hline
\hline
32 &  Easy\-ABC & 4 & \ra xml2ly \ra\ Elysium \ra\ & 2 & LilyPond & 5 &  &  \\
\hline
33 &  Easy\-ABC & 4 & \ra mxml2ly \ra\ Elysium \ra\ & 2 & LilyPond & 5 &  &  \\
\hline
34 &  Easy\-ABC & 4 & \ra xml2ly \ra\ Frescobaldi \ra\ & 2 & LilyPond & 5 &  &  \\
\hline
35 &  Easy\-ABC & 4 & \ra mxml2ly \ra\ Frescobaldi \ra\ & 2 & LilyPond & 5 &  &  \\
\hline
36 &  Rosegarden & 3 & \ra xml2ly \ra\ Elysium \ra\ & 2 & LilyPond & 5 & & \\
\hline
37 &  Rosegarden & 3 & \ra mxml2ly \ra\ Elysium \ra\ & 2 & LilyPond & 5 & & \\
\hline
38 &  Rosegarden & 3 & \ra xml2ly \ra\ Frescobaldi \ra\ & 2 & LilyPond & 5 & & \\
\hline
39 &  Rosegarden & 3 & \ra mxml2ly \ra\ Frescobaldi \ra\ & 2 & LilyPond & 5 & & \\
\hline
40 &  MuseScore & 5 & \ra xml2ly \ra\ Elysium \ra\ & 2 & LilyPond & 5 & & \\
\hline
41 &  MuseScore & 5 & \ra mxml2ly \ra\ Elysium \ra\ & 2 & LilyPond & 5 & & \\
\hline
42 &  MuseScore & 5 & \ra xml2ly \ra\ Frescobaldi \ra\ & 2 & LilyPond & 5 &  &  \\
\hline
43 &  MuseScore & 5 & \ra mxml2ly \ra\ Frescobaldi \ra\ & 2 & LilyPond & 5 &  &  \\
\hline
44 & Denemo & 3 & \ra xml2ly \ra\ Elysium \ra\ & 2 & LilyPond & 5 &  &  \\
\hline
45 & Denemo & 3 & \ra mxml2ly \ra\ Elysium \ra\ & 2 & LilyPond & 5 &  &  \\
\hline
46 & Denemo & 3 & \ra xml2ly \ra\ Frescobaldi \ra\ & 2 & LilyPond & 5 &  & \\
\hline
47 & Denemo & 3 & \ra mxml2ly \ra\ Frescobaldi \ra\ & 2 & LilyPond & 5 &  & \\
\hline
48 & Canorus & 3 & \ra xml2ly \ra\ Elysium \ra\ & 2 & LilyPond & 5 &  &  \\
\hline
49 & Canorus & 3 & \ra mxml2ly \ra\ Elysium \ra\ & 2 & LilyPond & 5 &  &  \\
\hline
50 & Canorus & 3 & \ra xml2ly \ra\ Frescobaldi \ra\ & 2 & LilyPond & 5 &  &  \\
\hline
51 & Canorus & 3 & \ra mxml2ly \ra\ Frescobaldi \ra\ & 2 & LilyPond & 5 &  &  \\
\hline
\hline
52 &  Text-Editor & 1 & \ra\ & 0 &  LilyPond & 5 & 1.0 & 0.6 \\
\hline
53 &  Elysium & 4 & \ra\ & 0 &  LilyPond & 5 & 1.0 & 0.9 \\
\hline
54 &  Frescobaldi & 5 & \ra\ & 0 &  LilyPond & 5 & 1.1 & 1.0 \\
\hline
\hline
\end{longtable}

\end{footnotesize}



% TODO nicht gegutachtete Wege: Fresco/ etc. nach Easyabc, da deaysy abc auch ein Texteditor


% this is only inserted to eject fault messages in texlipse
%\bibliography{../bib/literature}
